
Orbiter is a command line driven system.
There is no graphical user interface.
This means that commands are typed into a terminal, and executed by the operating system.
In this mode of operation, Orbiter is just like any other program installed on the computer. 
This also means that Orbiter 
can be mixed with other applications, using files to share data between the processes. 

\bigskip


The command line is entered into an application that is called {\em Terminal} (or {\em SuperShell} in Windows).
Orbiter is called from the command line, and command options are given to instruct Orbiter what to do.
The process that calls orbiter is the shell. There are different types of shells, but they all 
provide the necessary interface to allow the user to start jobs and maintain files. 
Shells can be programmed by means of shell scripts. Programming by mans of shell scripts is called scripting. 
Orbiter can be programmed using shell scripting. 


\bigskip


Make is a software tool that allows to execute small command sequences. 
Commands can trigger other commands, based on dependencies. The dependencies are tied to the existence of files, 
which can be named. If the file does not exist or is outdate, make triggers a command whose purpose is to create that file.
The dependencies are called rules. 
The makefile consists of rules and command sequences to be invoked based on the dependencies 
and the commands issued.
Many commands can be collected in one makefile. 
Because of the convenience offered by makefiles, 
this user's guide 
will rely on a the make / makefile tool.
The makefile associated with this user's guide is listed in full in Section~\ref{sec:makefile}.
It would also be possible to define shell scripts for each of the commands.


\bigskip

Make also allows to use variables, which are used by means of text substitution.
A variable 
is defined as 
\begin{verbatim}
A="I am a variable"
\end{verbatim}
and used anywhere later using the 
\begin{verbatim}
$(A)
\end{verbatim}
syntax. 
Rules are defined using the following syntax
\begin{verbatim}
Label:
  Do something
\end{verbatim}
Here, label is the name of the rule, and \verb'Do something' is the code that is exacuted whenever make is called 
with the given label in the command line. For instance
\begin{verbatim}
make Label
\end{verbatim}
will execute \verb'Do something'. 
The shell will take the command and peel off the first word, which is \verb'Do'. It will then search the 
system for a command called \verb'Do'.  Of course, this will result in an error because there is no command called \verb'Do'.
The remaining piece of the command line, i.e. \verb'something' is considered as an argument to the command. 
For instance, suppose we have a orbiter command with several options, say 
\begin{verbatim}
orbiter.out -v 3 -define F -finite_field -q 16 -end \
  -with F -do -finite_field_activity -cheat_sheet_GF -end
\end{verbatim}
The purpose of this command is to produce a file called 
$$
\verb'GF_16.tex'
$$ 
which can then be processed through latex to give the report. 
Observe that the command is quite long, and stretches over two lines. 
The backslash at the end of the first line indicates that the command continues on to the next line.
Using make, we can assign a label to this command. 
Suppose we want to call this command \verb'F_16'.
We can create a makefile like this:

\sourcecode{F_16}

With this file present, type the terminal command \verb'make F_16' to 
execute the two line Orbiter command.
Windows users can use \verb'SuperShell'.
The program \verb'make' will look for the file \verb'makefile' in the current directory.
Once found, it will search for the label \verb'F_16' in it and execute the commands beneath it.
The given commands will invoke Orbiter and produce the \verb'GF_16.tex' containing the desired report.
If we wanted to do some other Orbiter command, 
we could edit the makefile. 
We would also have a sequence of commands listed in the same target. 
In this case, makefile will process these commands one after the other.


\bigskip

Makefiles are somewhat picky when it comes to whitespacing.
The command sequence needs to be indented using tab symbols. 
Leading spaces will cause error messages. 
Also, there should be no whitespace after the trailing backslash symbol.
Some editors can display whitespace characters. This can be helpful when 
working with makefiles.



\bigskip

A sample makefile with all of the commands discussed in this user's guide 
is distributed with Orbiter (in the examples directory). 
The file is reproduced in Section~\ref{sec:makefile}. 
It is advised to copy the example makefile from the orbiter tree to a location outside the orbiter 
distribution directory (otherwise, git update will cause error messages). 
It is also fine to create a new custom makefile, considering the remarks about  \verb'ORBITER_PATH' below.



\bigskip


One difficulty in installing Orbiter is the path of installation.
In the sample makefile, there is a makefile variable called 
\verb'ORBITER_PATH' which contains the path to the orbiter executable \verb'orbiter.exe'.
Depending on the local installation of orbiter, 
the makefile variable needs to be changed accordingly. 
The actual command to run the \verb'F_16' example is as follows:
\begin{verbatim}
F_16:
  $(ORBITER_PATH)orbiter.out -v 3 -define F -finite_field -q 16 -end \
    -with F -do -finite_field_activity -cheat_sheet_GF -end
\end{verbatim}
The orbiter installation directory \verb'orbiter' 
and a second directory called \verb'work' should be next to each other. 
The orbiter example makefile should be copied into the \verb'work' directory. 
The top of the file should contain the line
\begin{verbatim}
MY_PATH=../orbiter
\end{verbatim}
This will set \verb'ORBITER_PATH' to point to the correct location of the orbiter executable. 
inside the \verb'work' directory, any of the commands listed in this guide will function correctly.
Another possibility is to install \verb'orbiter.out' in a central location, where it can be found by the shell. 
In this case, we should change the line 
\begin{verbatim}
ORBITER_PATH=$(MY_PATH)/src/apps/orbiter/
\end{verbatim}
to
\begin{verbatim}
ORBITER_PATH=
\end{verbatim}
in the makefile.
